% !TeX root = ../../elegantbook-cn.tex
% 所有题目都未审
\begin{example}
（2026东城一模）已知函数 $f(x)$ 的定义域为 $\left(\dfrac12,+\infty\right)$，$f(1)=1$，函数 $g(x)=x\ln(2x-1)$．对任意 $t\in\left(\dfrac12,+\infty\right)$，曲线 $y=f(x)$ 在点 $A(t,f(t))$ 处的切线方程为 $y-f(t)=g(t)(x-t)$．\\
（1）求 $f(x)$ 的最小值；\\
（2）讨论 $g(x)$ 的单调性；\\
（3）已知 $f(3)<3\ln5+1$，求过点 $(2,1)$ 且与曲线 $y=f(x)$ 相切的直线的条数．
\end{example}

\begin{example}
（2025东城一模）设函数 $f(x)=(x-2)\mathrm{e}^x+ax+b$，曲线 $y=f(x)$ 在 $A(0,f(0))$ 处的切线方程为 $y=0$．\\
（1）求 $a,b$ 的值；\\
（2）求不等式 $f(x)\geqslant0$ 的解集；\\
（3）已知 $B(t,f(t))$，其中 $t>2$，直线 $AB$ 的方程为 $y=g(x)$．若 $x_1,x_2\in(0,t)$，且 $f(x_1)=g(x_2)$，求证：$x_1>x_2$．
\end{example}

\begin{example}
（2024东城一模）已知函数 $f(x)=x\ln(x-1)$．\\
（1）求曲线 $y=f(x)$ 在 $x=2$ 处的切线方程；\\
（2）设 $g(x)=f'(x)$，求函数 $g(x)$ 的最小值；\\
（3）若 $\dfrac{f(x)}{x-a}>2$，求实数 $a$ 的值．
\end{example}

\begin{example}
（2023东城一模）已知函数 $f(x)=ax^2-x\ln x$．\\
（1）当 $a=0$ 时，求 $f(x)$ 的单调递增区间；\\
（2）设直线 $l$ 为曲线 $y=f(x)$ 的切线，当 $a\geqslant\dfrac{\mathrm{e}}2$ 时，记直线 $l$ 的斜率的最小值为 $g(a)$，求 $g(a)$ 的最小值；\\
（3）当 $a>0$ 时，设
\[
M=\left\{y\mid y=f'(x),\ x\in\left(\dfrac1{2a},\dfrac3{4a}\right)\right\},\quad
N=\left\{y\mid y=f'(x),\ x\in\left(\dfrac1{4a},\dfrac1{2a}\right)\right\},
\]
求证：$M\nsubseteq N$．
\end{example}

\begin{example}
（2022东城一模）已知函数 $f(x)=\dfrac{x-a}{x^2-1}$．\\
（1）若曲线 $y=f(x)$ 在点 $(2,f(2))$ 处的切线斜率为 $-1$，求 $a$ 的值；\\
（2）若 $f(x)$ 在 $(1,+\infty)$ 上有最大值，求 $a$ 的取值范围．
\end{example}

\begin{example}
（2021东城一模）已知函数 $f(x)=x^3-ax^2-a^2x+1$，其中 $a>0$．\\
（1）当 $a=1$ 时，求 $f(x)$ 的单调区间；\\
（2）若曲线 $y=f(x)$ 在点 $(-a,f(-a))$ 处的切线与 $y$ 轴的交点为 $(0,m)$，求 $m+\dfrac1a$ 的最小值．
\end{example}

\begin{example}
（2026西城一模）已知函数 $f(x)=x^2\ln x$，直线 $l$ 为曲线 $y=f(x)$ 在点 $(x_0,f(x_0))$ 处的切线，$O$ 为坐标原点．\\
（1）求函数 $f(x)$ 的单调区间；\\
（2）设 $g(x)=f(x)+x$，研究函数 $g(x)$ 的零点个数；\\
（3）若 $l$ 与 $x$ 轴、$y$ 轴分别交于点 $A,B$，且 $\triangle OAB$ 为等腰直角三角形，求 $\triangle OAB$ 的面积．
\end{example}

\begin{example}
（2025西城一模）已知函数 $f(x)=(x+a)\mathrm{e}^{ax}$，其中 $a\in\mathbb{R}$．\\
（1）若曲线 $y=f(x)$ 在点 $(0,f(0))$ 处的切线斜率为 $2$，求 $a$ 的值；\\
（2）求函数 $f(x)$ 的单调区间；\\
（3）设函数 $f(x)$ 在区间 $[-1,2]$ 上的最大值和最小值分别为 $M(a),N(a)$，求使得不等式 $M(a)N(a)\geqslant(2a-1)\mathrm{e}^a+\mathrm{e}^2$ 成立的 $a$ 的最小值．
\end{example}

\begin{example}
（2024西城一模）已知函数 $f(x)=x+\ln(ax)+\dfrac1a x\mathrm{e}^x$．\\
（1）当 $a=1$ 时，求曲线 $y=f(x)$ 在点 $(1,f(1))$ 处切线的斜率；\\
（2）当 $a=-1$ 时，讨论 $f(x)$ 的单调性；\\
（3）若集合 $\{x\mid f(x)\geqslant-1\}$ 有且只有一个元素，求 $a$ 的值．
\end{example}

\begin{example}
（2023西城一模）已知函数 $f(x)=\mathrm{e}^x-\cos x$．\\
（1）求曲线 $y=f(x)$ 在点 $(0,f(0))$ 处的切线方程；\\
（2）设 $g(x)=xf'(x)-f(x)$，证明：$g(x)$ 在 $(0,+\infty)$ 上单调递增；\\
（3）判断 $3f\left(\dfrac13\right)$ 与 $4f\left(\dfrac14\right)$ 的大小关系，并加以证明．
\end{example}

\begin{example}
（2022西城一模）已知函数 $f(x)=\dfrac{ax}{\mathrm{e}^x+a}-1$，$a\neq0$．\\
（1）当 $a=1$ 时，\\
①求曲线 $y=f(x)$ 在 $x=0$ 处的切线方程；\\
②求证：$f(x)$ 在 $(0,+\infty)$ 上有唯一极大值点；\\
（2）若 $f(x)$ 没有零点，求 $a$ 的取值范围．
\end{example}

\begin{example}
（2021西城一模）已知函数$f(x) = \mathrm{e}^x(\ln x - a)$.\\
（1）若$a = 1$，求曲线$y = f(x)$在点$(1, f(1))$处的切线方程；\\
（2）若$a > 1$，求证：函数$f(x)$存在极小值；\\
（3）若对任意的实数$x \in [1, +\infty)$，$f(x) \geq -1$恒成立，求实数$a$的取值范围.
\end{example}

\begin{example}
（2026朝阳一模）已知函数 $f(x)=x-x\ln x-k$，$g(x)=x\mathrm{e}^{1-x}-k$，其中 $k\in\mathbb{R}$．\\
（1）求 $f(x)$ 的最大值；\\
（2）若 $g(x)$ 在区间 $(0,1)$ 上有且只有一个零点 $m$，证明：$f(x)$ 在区间 $(1,+\infty)$ 上有且只有一个零点 $n$，且 $m\mathrm{e}^{n-1}=1$；\\
（3）对于（2）中的 $n$，证明：$g(n)>0$．
\end{example}

\begin{example}
（2025朝阳一模）已知函数 $f(x)=a\ln x-\dfrac{x-1}{x+1}$（$a\in\mathbb{R}$）．\\
（1）当 $a=1$ 时，求曲线 $y=f(x)$ 在点 $(1,f(1))$ 处的切线方程；\\
（2）若 $a\geqslant\dfrac12$，求证：当 $x\geqslant1$ 时，$f(x)\geqslant0$；\\
（3）若函数 $f(x)$ 有 $3$ 个不同的零点，求 $a$ 的取值范围．
\end{example}

\begin{example}
（2024朝阳一模）已知函数 $f(x)=(1-ax)\mathrm{e}^x$（$a\in\mathbb{R}$）．\\
（1）讨论 $f(x)$ 的单调性；\\
（2）若关于 $x$ 的不等式 $f(x)>a(1-x)$ 无整数解，求 $a$ 的取值范围．
\end{example}

\begin{example}
（2023朝阳一模）已知函数 $f(x)=\mathrm{e}^{2x}-ax-1$（$a\in\mathbb{R}$）．\\
（1）求 $f(x)$ 的单调区间；\\
（2）若 $f(x)>0$ 对 $x\in(0,+\infty)$ 恒成立，求 $a$ 的取值范围；\\
（3）证明：若 $f(x)$ 在区间 $(0,+\infty)$ 上存在唯一零点 $x_0$，则 $x_0<a-2$．
\end{example}

\begin{example}
（2022朝阳一模）已知 $f(x)=x-a\mathrm{e}^x$，$a\in\mathbb{R}$．\\
（1）若曲线 $y=f(x)$ 在点 $(1,f(1))$ 处的切线与 $x$ 轴重合，求 $a$ 的值；\\
（2）若函数 $f(x)$ 在区间 $(1,+\infty)$ 上存在极值，求 $a$ 的取值范围；\\
（3）设 $g(x)=f(2-x)$，在（2）的条件下，试判断函数 $g(x)$ 在区间 $(1,+\infty)$ 上的单调性，并说明理由．
\end{example}

\begin{example}
（2021朝阳一模）已知函数 $f(x)=(ax-1)\mathrm{e}^x$（$a\in\mathbb{R}$）．\\
（1）讨论 $f(x)$ 的单调性；\\
（2）若直线 $y=ax+a$ 与曲线 $y=f(x)$ 相切，求证：$-1<a<-\dfrac23$．
\end{example}

\begin{example}
（2026丰台一模）已知函数 $f(x)=\dfrac{x^3+a}{\mathrm{e}^{ax}}$．\\
（1）当 $a=1$ 时，求曲线 $y=f(x)$ 在 $(0,f(0))$ 处的切线方程；\\
（2）若 $a\geqslant\sqrt2$，证明：当 $x\in[0,+\infty)$ 时，$f(x)\leqslant a$；\\
（3）若存在 $x_0\in[0,+\infty)$，使 $f(x_0)<-a^2x_0+a$ 成立，求实数 $a$ 的取值范围．
\end{example}

\begin{example}
（2025丰台一模）已知函数 $f(x)=\ln x-\dfrac ax$，直线 $l$ 是曲线 $y=f(x)$ 在点 $(t,f(t))$（$t>0$）处的切线．\\
（1）当 $a=0$，$t=\mathrm{e}$（$\mathrm{e}$ 为自然对数的底数）时，求 $l$ 的方程；\\
（2）若存在 $l$ 经过点 $(0,0)$，求实数 $a$ 的取值范围；\\
（3）当 $a=-1$ 时，设点 $A(t,f(t))$（$t>0$），$O(0,0)$，$B$ 为 $l$ 与 $y$ 轴的交点，$S_{\triangle AOB}$ 表示 $\triangle AOB$ 的面积，求 $S_{\triangle AOB}$ 的最小值．
\end{example}

\begin{example}
（2024丰台一模）已知函数 $f(x)=\mathrm{e}^x+\ln(x+1)-x$，曲线 $C:y=f(x)$ 在点 $(x_0,f(x_0))$ 处的切线为 $l:y=g(x)$，记 $h(x)=f(x)-g(x)$．\\
（1）当 $x_0=0$ 时，求切线 $l$ 的方程；\\
（2）在（1）的条件下，求函数 $h(x)$ 的零点并证明对 $x$ 有 $h(x)\geqslant0$；\\
（3）当 $x_0\neq0$ 时，直接写出函数 $h(x)$ 的零点个数（结论不要求证明）．
\end{example}

\begin{example}
（2023丰台一模）已知函数 $f(x)=x+\dfrac a{\mathrm{e}^x}$（$a>0$）．\\
（1）求函数 $f(x)$ 的极值；\\
（2）若函数 $f(x)$ 有两个不相等的零点 $x_1,x_2$，\\
（i）求 $a$ 的取值范围；\\
（ii）证明：$x_1+x_2>2\ln a$．
\end{example}

\begin{example}
（2022丰台一模）已知函数 $f(x)=x\sqrt{a-x}$．\\
（1）当 $a=1$ 时，求曲线 $y=f(x)$ 的斜率为 $1$ 的切线方程；\\
（2）若函数 $g(x)=f(x)-\dfrac{2a}{3}$ 恰有两个不同的零点，求 $a$ 的取值范围．
\end{example}

\begin{example}
（2021丰台一模）已知函数 $f(x)=x^3-3x^2+b$（$b\in\mathbb{R}$）．\\
（1）当 $b=1$ 时，求曲线 $y=f(x)$ 在点 $(1,f(1))$ 处的切线方程；\\
（2）若函数 $f(x)$ 存在三个零点，分别记为 $x_1,x_2,x_3$（$x_1<x_2<x_3$），\\
（i）求 $b$ 的取值范围；\\
（ii）证明：$x_1+x_2>0$．
\end{example}

\begin{example}
（2026海淀一模）设函数 $f(x)=x^2-\ln(ax+1)+ax$（$a\neq0$）．\\
（1）当 $a=1$ 时，求证：直线 $y=0$ 是曲线 $y=f(x)$ 的切线；\\
（2）求 $f(x)$ 的单调区间；\\
（3）判断函数 $g(x)=xf(x)$ 是否存在极值．如果存在，求出所有的极值；如果不存在，说明理由．
\end{example}

\begin{example}
（2025海淀一模）已知函数 $f(x)=-\sin x+kx$．\\
（1）若曲线 $y=f(x)$ 在点 $(0,f(0))$ 处的切线为 $y=-\dfrac32x$，求 $k$ 的值；\\
（2）若 $f(x)$ 为 $\mathbb{R}$ 上的单调函数，求 $k$ 的取值范围；\\
（3）若函数 $g(x)=\dfrac{k}{6}x^3+x+\sin x$，求证：$k$ 可以取无数个值，使得每一个 $k$ 的取值对应的 $g(x)$ 都恰有三个不同的零点．
\end{example}

\begin{example}
（2024海淀一模）已知函数 $f(x)=x\mathrm{e}^{a-\frac12x}$．\\
（1）求 $f(x)$ 的单调区间；\\
（2）若函数 $g(x)=\left|f(x)+\mathrm{e}^{-2a}\right|$（$x\in(0,+\infty)$）存在最大值，求 $a$ 的取值范围．
\end{example}

\begin{example}
（2023海淀一模）已知函数 $f(x)=\mathrm{e}^{ax}-x$．\\
（1）当 $a=1$ 时，求曲线 $y=f(x)$ 在点 $(0,f(0))$ 处的切线方程；\\
（2）求 $f(x)$ 的单调区间；\\
（3）若存在 $x_1,x_2\in[-1,1]$，使得 $f(x_1)f(x_2)\geqslant9$，求 $a$ 的取值范围．
\end{example}

\begin{example}
（2022海淀一模）已知函数 $f(x)=\mathrm{e}^x(ax^2-x+1)$．\\
（1）求曲线 $y=f(x)$ 在点 $(0,f(0))$ 处的切线方程；\\
（2）若函数 $f(x)$ 在 $x=0$ 处取得极大值，求 $a$ 的取值范围；\\
（3）若函数 $f(x)$ 存在最小值，直接写出 $a$ 的取值范围．
\end{example}

\begin{example}
（2021海淀一模）已知函数 $f(x)=x\sin x$．\\
（1）判断函数 $f(x)$ 在区间 $\left(0,\dfrac\pi2\right)$ 上的单调性，并说明理由；\\
（2）求证：函数 $f(x)$ 在 $\left(\dfrac\pi2,\pi\right)$ 内有且只有一个极值点；\\
（3）求函数 $g(x)=\dfrac{f(x)+2}{\ln x}$ 在区间 $(1,\pi]$ 上的最小值．
\end{example}
